\documentclass[12pt]{article}
\usepackage[utf8]{inputenc}
\usepackage[T1]{fontenc}
\usepackage{hyperref}
\title{Yugh Language Resources for K-12 Curriculum Development}
\author{Compiled by Claude Sonnet 5.5 (claude-sonnet-5-5) for Andr\'as Kornai}
\date{}
\begin{document}
\maketitle

\section{Introduction}
Yugh (also Sym Ket, Southern Ket; ISO 639-3 yuu) is a Yeniseian language of central Siberia, closely related to Ket. It is extinct. Fluent speakers were gone by the early 1990s, and Heinrich Werner documented the last speakers in 1961--1971. The family label ``Bookkeeping'' in the request appears to be a catalogue placeholder; sources found describe Yugh as Yeniseian. No school curriculum in Yugh exists, and building one would be a revitalization or reconstruction project based on archival documentation only. Only the resources below were seen in search results.

\section{Primary Educational Resources}
No primers, textbooks or school materials for Yugh were found in search results.

\section{Linguistic and Reference Materials}
\begin{itemize}
\item Werner and Vajda, \emph{Dictionary of the Yugh language} (Lincom, 2012), based on Werner's fieldwork with the last speakers. It includes a description of the language and an extensive vocabulary. Search results referred to it, but no publisher page was among them. Related page: \url{https://works.bepress.com/edward_vajda/74}
\item Wikipedia overview of the Yugh language: \url{https://en.wikipedia.org/wiki/Yugh_language}
\item Glottolog record for the language, which lists older sources such as Castr\'en (1858) and Werner (1997), \emph{Das Jugische (Sym-Ketische)}: \url{https://glottolog.org/resource/languoid/iso/yug}
\item Endangered Languages Project entries: \url{https://endangeredlanguages.com/node/95312} and \url{https://endangeredlanguages.com/node/94476}
\end{itemize}

\section{Religious and Community Resources}
No scripture translations or community organizations specific to Yugh were found.

\section{Digital and Audio Resources}
No audio, video or corpus resources were found in the search results.

\section{Curriculum Development Context}
Yugh is extinct, so a K-12 curriculum would have to be a heritage or reconstruction programme. The Werner--Vajda dictionary and the Glottolog-listed grammars would be the core sources. The closely related living Ket language, which is itself endangered, could serve as a comparative bridge. Little else is available.
\end{document}